\documentclass[11pt,letterpaper]{article}
\usepackage[margin=1.05in]{geometry}
\usepackage{amsmath,amssymb,amsthm,xcolor}
\usepackage[colorlinks=true,linkcolor=teal!60!black,citecolor=teal!60!black,urlcolor=teal!60!black]{hyperref}
\hypersetup{pdftitle={How many integer roots can a short arithmetic program have?}}
\newtheorem{theorem}{Theorem}
\newtheorem{lemma}{Lemma}
\newcommand{\Z}{\mathbb Z}
\newcommand{\Q}{\mathbb Q}
\setlength{\parindent}{0pt}
\setlength{\parskip}{5pt plus 1pt minus 1pt}
\raggedbottom
\title{\bfseries How many integer roots can a short\\arithmetic program have?}
\author{A research note on the Shub--Smale conjecture}
\date{November 2026}
\begin{document}
\maketitle
\begin{abstract}
How many integer roots can a polynomial have if it takes only \(s\)
arithmetic operations to compute? The Shub--Smale conjecture predicts a
polynomial bound. We derive an exponential upper bound smaller than the
degree bound \(2^s\) by a factor of order \(\sqrt{s}\,2^{\sqrt{2s}}\).
The proof combines a real fewnomial theorem with the coefficient growth
forced by distinct integer roots. A classical construction gives a nearly
quadratic lower bound. Tracking gate order saves an additional factor of
\(\sqrt{s}\) over the same argument using only gate counts.
\end{abstract}

\section{Introduction}\label{sec:intro}
A straight-line program starts with \(1,X\). Each gate adds, subtracts or
multiplies two earlier values, and a value can be reused. Let \(\tau(f)\)
be the fewest gates computing \(f\in\Z[X]\). Let \(Z(s)\) be the maximum
number of \emph{distinct integer roots} of a nonzero \(f\) with \(\tau(f)\le s\).
Each multiplication can at most double the degree, giving \(Z(s)\le 2^s\).
Shub and Smale ask whether \(Z(s)\le (1+s)^{O(1)}\) \cite{SS,Smale}.

\begin{theorem}\label{thm:bounds}
For absolute constants \(c,C>0\) and all sufficiently large \(s\),
\[
 c\,\frac{s^2}{(\log s)^2(\log\log s)^2}
 \ \le\ Z(s)\ \le\
 C\,\frac{2^{s-\sqrt{2s}}}{\sqrt{s}} .
\]
All logarithms are to base two.
\end{theorem}

The lower estimate follows from published algorithms. The upper estimate
is a deduction here from published real-root bounds and the order of the
gates. Its literature priority is undetermined. Section~\ref{sec:gain} compares an earlier arithmetic estimate. Both bounds allow
\(Z(s)=O(s^2)\).

Section~\ref{sec:lower} constructs many roots using fast evaluation of a
product. Section~\ref{sec:upper} runs the argument in the other direction.
Given many integer roots, how large must the coefficients be, and how
many additive gates must the program use? A single choice of where to
split the program lets us combine these two constraints.

\clearpage
\section{Prior work and the lower bound}\label{sec:lower}
Smale credits Strassen with showing that the exponent in any polynomial
bound for \(Z(s)\) must be at least two \cite[Problem 4]{Smale}.
Koiran, Portier and Rojas describe the falling-factorial argument using
fast multipoint evaluation \cite[after Theorem 1.14]{KPR}. Put
\[
 P_N(X)=(X-1)(X-2)\cdots(X-N).
\]
A short program for any nonzero integer multiple of \(P_N\) gives a
\emph{lower} bound for \(Z(s)\), because the output vanishes at all
\(N\) integers from 1 through \(N\).

\paragraph{Computing the product in blocks.}
For \(N=h^2\), let \(B(T)=\prod_{j=1}^h(T-j)\). The desired product is
\(\prod_{i=0}^{h-1} B(X-ih)\). There are \(h\) blocks, and the inputs to
\(B\) form an arithmetic progression. Fast evaluation shares work across
all the blocks. To keep the logarithmic factors precise, use the recurrence
algorithm of Bostan, Gaudry and Schost \cite[Theorem 14]{BGS}. Over the
ring \(\Q[X]\), its scalar recurrence
\[
 U_0=1,\qquad U_i=(X-i)U_{i-1}
\]
gives \(U_N=P_N\) in \(O(M(\lceil\sqrt N\rceil))\) ring operations.
Here \(M(h)\) counts the ring operations to multiply two polynomials of
degree less than \(h\) in an auxiliary variable.
Cantor and Kaltofen give \(M(h)=O(h\log h\log\log h)\) for large \(h\)
\cite{CK}. Bostan and Yurkevich give a related algorithm for blocks
whose inputs form a geometric progression \cite{BY}.

\begin{lemma}[Clearing scalar denominators]\label{lem:scalar}
Suppose a program starts with \(1,X\) and uses arithmetic and divisions
by nonzero rational scalars generated by the program. A division-free
program of length at most a constant times its length computes a nonzero
integer multiple of its output whenever that output is in \(\Z[X]\).
\end{lemma}
\begin{proof}
Store each register as a pair \((A,d)\) meaning \(A/d\), where \(A\in\Z[X]\)
and \(0\ne d\in\Z\). Use the usual common-denominator rule for addition
and subtraction, and multiply the two coordinates separately for products.
Dividing \(A/d\) by a scalar \(u/v\) gives the pair \((Av,du)\). Each update
takes a bounded number of gates. Return the final numerator.
\end{proof}

The scalar inverses needed in the recurrence algorithm are inverses of
explicit integers \(D(1,2^t,2^t)\), with \(0\le t\le\lceil\log_4 N\rceil\).
Lemma 6 of \cite{BGS} constructs each denominator in \(O(2^t)\) operations.
The powers of two are obtained by doubling, so scalar preparation costs
\(O(\sqrt N)\) in total. Cantor--Kaltofen multiplication can use its
integer bilinear algorithm, which realizes its integer scaling by repeated additions and subtractions
\cite{CK}. Lemma~\ref{lem:scalar} thus gives the same roots within
\(O(h\log h\log\log h)\) gates. Choosing \(h\) of order
\(s/(\log s\log\log s)\) with a sufficiently small constant proves the
lower half of Theorem~\ref{thm:bounds}.

\clearpage
\section{Bounding integer roots from above}\label{sec:upper}
Fix a program with \(a\) additive gates (additions and subtractions) and
\(m\) multiplications, with \(a+m\le s\). Write \(R\) for the number of
distinct integer roots of its nonzero output \(f\), and \(L=\log R\).
We only need the case of large \(R\).
Risler proved that \(R\le C_0^{a^2}\) by bounding all real roots
\cite{Risler}. Together with \(R\le 2^m\), this already implies
\(s\ge L+\Omega(\sqrt L)\). To sharpen the saving we count \emph{when}
the additions took place.

\subsection{How large must the coefficients get?}
Stop the program after any gate and subtract the number of multiplications
so far from the number of additive gates so far. Let \(d\) be the
maximum of this difference, including the empty prefix. Thus \(0\le d\le a\).
Early additions matter because later multiplications can repeatedly square
what they produced.

\begin{lemma}[Height]\label{lem:height}
With \(L=\log R\) and \(d\) the maximum prefix difference just defined,
\[
 s\ \ge\ L+\log L+a-\log(d+1)-O(1).
 \tag{H}\label{eq:H}
\]
\end{lemma}
\begin{proof}
Write \(\|g\|_1\) for the sum of the absolute values of the coefficients of
\(g\). Let \(H\) be the log of the largest such norm in any register so far,
including the initial \(1\). At the start \(H=0\). An addition raises \(H\)
by at most one. A multiplication at most doubles it. If \(m_i\) products
precede additive gate \(i\), then
\[
 2^{-m}\log\|f\|_1\le \sum_i2^{-m_i}\le d+1.
\]
Indeed, at most \(j+d\) additions have \(m_i\le j\), by the prefix
definition. Writing \(2^{-m_i}=\sum_{j\ge m_i}2^{-j-1}\) bounds the sum
by \(\sum_{j\ge0}(j+d)2^{-j-1}=1+d\).

For \(f(X)=c\prod_j(X-z_j)\), its Mahler measure is
\(\mathcal M(f)=|c|\prod_j\max(1,|z_j|)\).
The geometric mean of \(|f|\) on the unit circle equals \(\mathcal M(f)\)
by Jensen's formula, so \(\mathcal M(f)\le\|f\|_1\).
Integrality also enters at this point. The leading coefficient satisfies \(|c|\ge1\).
At least half of \(R\) distinct integers have absolute value at least
\(R/8\) for large \(R\). Hence \(\log\|f\|_1\ge (R/2)\log(R/8)\).
Comparing the two bounds and taking one more logarithm gives
\(m\ge L+\log L-\log(d+1)-O(1)\). Add \(a\).
\end{proof}

\clearpage
\subsection{How many real roots can the gates describe?}
A fewnomial system is a system of polynomial equations with few distinct
monomials across all equations. Its solutions will encode the values of
registers at a root. Choose a prefix attaining the maximum difference
\(d\), the maximum of additive gates minus products in a prefix.
Here \(a\) counts all additive gates. Put \(b=a-d\). A larger lead \(d\)
means fewer extra monomials are needed to describe the remaining gates.

\begin{lemma}[Real roots]\label{lem:real}
For \(L=\log R\), where \(R\) counts integer roots of the program output,
\[
 L\ \le\ \tfrac12 b^2+b\log(2a+1)+O(a+1).
 \tag{F}\label{eq:F}
\]
\end{lemma}
\begin{proof}
Let the chosen prefix contain \(u\) additions or subtractions and \(v\)
multiplications, so \(u-v=d\). Keep one variable for \(X\) and one for
each nonconstant gate output in the prefix. If this gives \(n_0\) variables,
the constant and the coordinates account for at most \(1+n_0\) monomials.
Each product introduces at most one more, for a total of at most \(1+n_0+v\).
Replace constant registers by their values throughout.

In the suffix, introduce a variable only for a nonconstant additive gate.
Every intervening product is a scalar times a monomial in existing
variables. On points with nonzero coordinates, divide each new defining
equation by one of its terms. This gives a Laurent polynomial (negative
exponents are allowed) using the shared constant and at most two further
monomials. Finally append \(f^2=\varepsilon\), for a positive scalar
\(\varepsilon\). In these variables \(f\) is a scalar times a monomial,
so this adds at most one monomial. With \(t\le a-u\) new additive
variables, the total dimension is \(n=n_0+t\le1+a+v\le 2a+1\).
Here the support is the set of distinct monomials across all equations. Its size is at most
\[
 (1+n_0+v)+2t+1=n+1+(v+t+1)\le n+1+(b+1).
\]

The final inequality uses \(v+t\le v+a-u=a-d=b\). This is where the
prefix chosen for the height bound also reduces the monomial count.

For small positive \(\varepsilon\), each integer root of \(f\) has two
nearby solutions of \(f^2=\varepsilon\), one on each side. Use disjoint
neighborhoods and choose \(\varepsilon\) outside the finitely many critical
values of \(f^2\) and the values of \(f^2\) at zeros of \(X\) or of any
nonconstant register. Each solution lifts to a point with nonzero
coordinates. The equations are nonsingular there because the register
equations define a triangular graph and the final level has nonzero
derivative along it. Dividing equations by nonvanishing monomials
preserves nonsingularity.

Bihan and Sottile \cite{BS} bound the number of nonsingular positive
solutions for \(n\) equations with \(n+1+k\) monomials by
\(K2^{\binom{k}{2}}n^k\), with \(K=(e^2+3)/4\).
Apply the theorem after each change of coordinate signs. A smaller support
also satisfies the estimate, since at a nonsingular root the exponent differences span \(n\) dimensions,
and the bound increases with the excess above \(n+1\).
There are at most \(2^n\) sign choices. Taking \(k=b+1\) therefore gives
\(2R\le 2^n K2^{\binom{k}{2}}n^k\). Taking logs and using \(b\le a\)
proves \eqref{eq:F}.
\end{proof}

\clearpage
\subsection{The gain from tracking gate order}\label{sec:gain}
Recall that \(L=\log R\), \(a\) counts additive gates, \(d\) is the largest
excess of additions over multiplications in a prefix, and \(b=a-d\).
The height bound \eqref{eq:H} and the fewnomial bound \eqref{eq:F} now
give the upper half of Theorem~\ref{thm:bounds}.

\begin{proof}[Proof of the upper bound]
We can assume \(L\ge s-2\sqrt{s}\), since smaller \(L\) already satisfies
the claimed estimate for large \(s\). The degree bound gives \(L\le m\le s-a\).
Thus \(a\le 2\sqrt{s}\), and \(\log L=\log s+O(1)\).
Substituting these estimates in \eqref{eq:F} gives
\[
 s-O(\sqrt{s})\le \tfrac12 b^2+\tfrac12 b\log s+O(\sqrt{s}).
\]
Completing the square and taking the positive square root yields
\[
 b\ge\sqrt{2s}-\tfrac12\log s-O(1).
\]
Indeed an \(O(\sqrt{s})\) change under this square root costs only \(O(1)\).
Now \eqref{eq:H}, with \(a=b+d\), gives
\[
 L\le s-b-\log L-\bigl(d-\log(d+1)\bigr)+O(1).
\]
The parenthesized term is nonnegative for integers \(d\ge0\).
Substitution yields \(L\le s-\sqrt{2s}-\tfrac12\log s+O(1)\).
Exponentiation gives the result.
\end{proof}

\paragraph{A comparison using only gate counts.}
One can use the empty prefix in the fewnomial construction. Its monomial
excess is at most \(a+1\), so the preceding square-completion gives the
same lower bound for \(a\) as it did for \(b\). Without using gate order,
the height calculation gives \(\log\|f\|_1\le a2^m\), hence
\(L\le s-a-\log L+\log a+O(1)\). Using \(a\le 2\sqrt{s}\) in the
large-\(L\) regime gives \(L\le s-\sqrt{2s}+O(1)\). Thus gate order
saves a further factor of \(\sqrt{s}\) over this particular deduction.

\paragraph{Scope of related arithmetic estimates.}
Rojas \cite{Rojas} states a rational-root estimate which would give
\(L=O(a\log(a+1))\). Together with \(L\le s-a\) this would imply the
stronger saving \(L\le s-\Omega(s/\log s)\), since for \(L\ge s/2\) the
first inequality forces \(a=\Omega(s/\log s)\). Its proof invokes a multivariate
\(p\)-adic root bound. Goldstein, Rojas, Stone and Vera \cite{GRSV}
describe the earlier explicit multivariate bounds as requiring tropical
genericity and treat a particular nongeneric class. We have not verified
the needed hypothesis for arbitrary circuit encodings or an independent
proof of the unrestricted corollary. The real argument here uses the
regularity supplied by the level \(f^2=\varepsilon\).

\section{Conclusion}
Let \(T(N)\) be the minimum number of gates needed to compute a nonzero
integer polynomial vanishing at \(1,\ldots,N\). The lower-bound construction gives
\(T(N)=O(\sqrt N\log N\log\log N)\). In the other direction, substitute
\(s=T(N)\) and \(L\ge\log N\) in the upper bound. Since \(s\ge\log N\),
\[
 T(N)\ge \log N+\sqrt{2\log N}+\tfrac12\log\log N-O(1).
\]
This gap persists even with the roots fixed in advance. For arbitrary
sets of integer roots, whether \(Z(s)=O(s^2)\) remains open.

\clearpage
\begin{thebibliography}{99}\small
\bibitem{SS} M. Shub and S. Smale.
\href{https://janos.cs.technion.ac.il/COURSES/238900-13/Papers-1/ShubSmale-pap97.pdf}{On the intractability of Hilbert's Nullstellensatz and an algebraic version of ``NP \(\ne\) P?''}
\emph{Duke Mathematical Journal} 81, 47--54, 1995.
\bibitem{Smale} S. Smale.
\href{https://www.fim.uni-passau.de/fileadmin/dokumente/fakultaeten/fim/lehrstuhl/muller/SmaleProblems1998.pdf}{Mathematical problems for the next century}.
\emph{Mathematical Intelligencer} 20, 7--15, 1998. See Problem 4.
\bibitem{KPR} P. Koiran, N. Portier and J. M. Rojas.
\href{https://arxiv.org/abs/1309.0486}{Counting tropically degenerate valuations and \(p\)-adic approaches to the hardness of the permanent}.
arXiv:1309.0486, 2013.
\bibitem{BGS} A. Bostan, P. Gaudry and É. Schost.
\href{https://specfun.inria.fr/bostan/publications/BoGaSc07.pdf}{Linear recurrences with polynomial coefficients and application to integer factorization and Cartier--Manin operator}.
\emph{SIAM Journal on Computing} 36, 1777--1806, 2007.
\bibitem{CK} D. G. Cantor and E. Kaltofen.
\href{https://users.cs.duke.edu/~elk27/bibliography/91/CaKa91.pdf}{On fast multiplication of polynomials over arbitrary algebras}.
\emph{Acta Informatica} 28, 693--701, 1991.
\bibitem{BY} A. Bostan and S. Yurkevich.
\href{https://arxiv.org/abs/2012.08656}{Fast computation of the \(N\)-th term of a \(q\)-holonomic sequence and applications}.
\emph{Journal of Symbolic Computation} 115, 96--123, 2023.
\bibitem{Risler} J.-J. Risler.
\href{https://doi.org/10.1137/0214014}{Additive complexity and zeros of real polynomials}.
\emph{SIAM Journal on Computing} 14, 178--183, 1985.
\bibitem{BS} F. Bihan and F. Sottile.
\href{https://arxiv.org/abs/math/0609544}{New fewnomial upper bounds from Gale dual polynomial systems}.
\emph{Moscow Mathematical Journal} 7, 387--407, 2007.
\bibitem{Rojas} J. M. Rojas.
\href{https://arxiv.org/abs/math/0201154}{Additive complexity and the roots of polynomials over number fields and \(p\)-adic fields}.
\emph{ANTS V}, 506--515, 2002.
\bibitem{GRSV} J. Goldstein, J. M. Rojas, H. Stone and A. Vera.
\href{https://doi.org/10.1145/3832046.3832061}{New explicit root counts for \(p\)-adic circuit systems}.
\emph{ACM Communications in Computer Algebra} 59(4), 134--137 (online 2026).
\end{thebibliography}
\end{document}
